
    \subsection{R-to-T trefoil closure bridge}
        \label{subsec:rt_trefoil_closure_bridge}

        \textbf{[SPECULATIVE BRIDGE].}
        The photon-like R-phase packet is interpreted as a phase-torsion pulse
        on a swirl-characteristic line, not as a pre-existing broken vortex tube.
        A Compton-scale torsion packet may close into an electron-candidate
        trefoil only when a local boundary condition supplies trapping, recoil,
        and topological compensation. The minimal bridge statement is
        \begin{align}
            \gamma_{\rm R}
            +
            \mathcal{B}_{\rm atom}
            \overset{\mathcal{K}_{\rm Kairos}}{\longrightarrow}
            e^-_{3_1}
            +
            \mathcal{R},
            \label{eq:r_phase_to_trefoil_electron}
        \end{align}
        with closure conditions
        \begin{align}
            \Delta\Theta
            &=
            \int_{s_1}^{s_2}\partial_s\theta(s,t)\,ds
            \approx \pi,\\
            E_{\rm trap} &\gtrsim m_e c^2,\\
            \Gamma_{\rm trap} &\approx \Gamma_0,\\
            d_{\rm eye} &\sim \mathcal{O}(a_{\rm core}).
            \label{eq:trefoil_closure_conditions}
        \end{align}
        The phase condition \(\Delta\Theta\approx\pi\) corresponds to the
        minimal twist-knot closure \(m=1\), hence to \(3_1\) in the convention of
        Section~\ref{subsec:twist_ladder_closure_convention}.

        \textbf{[CRITICAL NOTE].}
        Equation~\eqref{eq:r_phase_to_trefoil_electron} is not a pure Euler
        reconnection process and is not a claim that a single free photon creates
        an isolated electron from vacuum. In a free process, charge, energy,
        momentum, and topological compensation require a conjugate sector or an
        external field reservoir. The canon-safe reading is closure of a trapped
        or incipient trefoil carrier under a boundary-supported R-to-T transition.

    \subsection{Photon torsion packets and R-to-T closure bookkeeping}
        \label{subsec:photon_torsion_r_to_t_closure}

        \textbf{[DERIVED SST CONSTRAINT].}
        The radiative photon sector is not represented as a free end of a vortex
        filament. A photon is represented as an extended R-phase torsion packet
        on a swirl-characteristic or transverse director field. Before trapping
        and closure, it has no protected closed-knot type:
        \begin{align}
            K(\gamma_{\rm R})\quad\text{is not a closed matter-knot invariant.}
        \end{align}

        \textbf{[SPECULATIVE SST BRIDGE].}
        Electron-trefoil formation from radiation is therefore admissible only
        as a boundary-supported R-to-T conversion, not as bulk reconnection of
        an already closed Euler vortex tube. The minimal bookkeeping reaction is
        \begin{align}
            \gamma_{\rm R}
            +
            \mathcal{B}_{\rm atom}
            \overset{\mathcal{K}_{\rm Kairos}}{\longrightarrow}
            e^-_{3_1}
            +
            \mathcal{R},
            \label{eq:photon_r_to_t_kairos_closure}
        \end{align}
        where \(\mathcal{B}_{\rm atom}\) supplies the trapping/boundary
        condition and \(\mathcal{R}\) carries recoil, momentum balance, and
        topological compensation.

        \textbf{[ORTHODOX CONSTRAINT / SST BOOKKEEPING].}
        A single free photon cannot be promoted to a lone charged electron in
        empty space without a reservoir for momentum and charge/topological
        compensation. Free matter creation must instead involve a conjugate
        sector or an external field reservoir, schematically
        \begin{align}
            \gamma+\gamma
            \longrightarrow
            3_1+\overline{3_1},
            \qquad
            \gamma+\mathcal{F}
            \longrightarrow
            3_1+\overline{3_1}+\mathcal{F}' .
            \label{eq:pair_creation_topological_bookkeeping}
        \end{align}

        The Compton threshold bookkeeping is
        \begin{align}
            \omega_c = \frac{\vnorm}{\rc},
            \qquad
            \hbar\omega_c=m_e c^2,
            \qquad
            E_{\rm pair}=2m_e c^2 .
            \label{eq:compton_pair_threshold_bookkeeping}
        \end{align}
        Thus \(\hbar\omega_c\) is the natural seed-closure scale for an
        already boundary-supported trefoil event, while \(2\hbar\omega_c\) is
        the natural pair-production scale for producing a trefoil and an
        anti-trefoil sector.
